% Standalone-Diagramm des Persistenz-/Verzeichnislayouts.
% Kompilieren:  pdflatex persistenzlayout.tex
% PNG:          pdftoppm -png -r 300 persistenzlayout.pdf persistenzlayout
\documentclass[border=12pt]{standalone}
\usepackage[utf8]{inputenc}
\usepackage[T1]{fontenc}
\usepackage[ngerman]{babel}
\usepackage{tikz}
\usetikzlibrary{positioning,arrows.meta}

\definecolor{cData}{HTML}{DCE7F5}
\definecolor{cDataB}{HTML}{5B7FB0}
\definecolor{cModel}{HTML}{DDEFDD}
\definecolor{cModelB}{HTML}{5A8F5A}
\definecolor{cStore}{HTML}{EFEFEF}
\definecolor{cStoreB}{HTML}{7A7A7A}
\definecolor{cRQ}{HTML}{7A7A7A}

\begin{document}
\begin{tikzpicture}[
  font=\small\sffamily,
  dir/.style={draw, rounded corners=2pt, align=left, minimum height=8mm,
              inner sep=6pt, line width=0.7pt, fill=cData, draw=cDataB,
              font=\fontsize{12.2}{14.4}\selectfont\bfseries\ttfamily},
  file/.style={draw, rounded corners=2pt, align=left, minimum height=7mm,
               inner sep=5pt, line width=0.5pt, fill=white, draw=cStoreB,
               font=\fontsize{10.8}{12.7}\selectfont\ttfamily},
  note/.style={font=\fontsize{10.8}{12.7}\selectfont\itshape\sffamily, text=cRQ, align=left, text width=46mm},
  edge/.style={draw=cStoreB, line width=0.5pt},
]

% --- Zweig models/ ----------------------------------------------------------
\node[dir] (models) {models/};
\node[file, below=6mm of models.south west, anchor=north west, xshift=8mm] (m1)
  {driving\_<source>.joblib};
\node[file, below=3mm of m1.south west, anchor=north west, xshift=0mm] (m2)
  {<source>\_forecaster\_<algo>.joblib};
\node[note, right=3mm of m1] {Driving-Classifier je Quelle};
\node[note, right=3mm of m2] {Forecaster je Quelle\,$\times$\,Algorithmus (rf/lgbm/lstm)};

% --- Zweig predictions/ -----------------------------------------------------
\node[dir, below=16mm of m2.south west, anchor=north west] (preds) {predictions/};
\node[dir, below=6mm of preds.south west, anchor=north west, xshift=6mm] (model)
  {<model>/};
\node[dir, below=4mm of model.south west, anchor=north west, xshift=6mm] (run)
  {<dataset>\_<N>d\_T<DD-MM-YYYY>/};
\node[file, below=4mm of run.south west, anchor=north west, xshift=6mm] (f1)
  {forecast.csv};
\node[file, below=3mm of f1.south west, anchor=north west] (f2)
  {*.png \textperiodcentered\ notes.txt};

\node[note, right=3mm of run] {eine Run-Dir = eine Simulation \textperiodcentered\ Helper \texttt{run\_output\_dir()}};
\node[note, right=3mm of f1] {Prognose (Trips, Quantile)};
\node[note, right=3mm of f2] {optionale Artefakte};

% --- Baumkanten -------------------------------------------------------------
\draw[edge] (models.south west)+(3mm,0) |- (m1.west);
\draw[edge] (models.south west)+(3mm,0) |- (m2.west);
\draw[edge] (preds.south west)+(3mm,0) |- (model.west);
\draw[edge] (model.south west)+(3mm,0) |- (run.west);
\draw[edge] (run.south west)+(3mm,0) |- (f1.west);
\draw[edge] (run.south west)+(3mm,0) |- (f2.west);

\end{tikzpicture}
\end{document}
